\documentclass[conference]{IEEEtran}

\usepackage{amsmath}
\usepackage{booktabs}
\usepackage{float}
\usepackage{graphicx}
\usepackage{microtype}
\usepackage{siunitx}
\usepackage[hidelinks]{hyperref}

\hypersetup{
  pdftitle={Validation-Guided Fusion of Cough and Speech for Participant-Disjoint COVID-19 Respiratory-Audio Classification},
  pdfauthor={Nishant Harkut and Ritu Tiwari},
  pdfkeywords={COVID-19, respiratory audio, multimodal fusion, participant-disjoint evaluation}
}

\graphicspath{{figures/}}
\sisetup{detect-all,round-mode=places,round-precision=3}

\title{Validation-Guided Fusion of Cough and Speech for Participant-Disjoint COVID-19 Respiratory-Audio Classification}

\author{
\IEEEauthorblockN{Nishant Harkut and Ritu Tiwari}
\IEEEauthorblockA{Department of Information Technology\\
Atal Bihari Vajpayee Indian Institute of Information Technology and Management\\
Gwalior, India}
}

\begin{document}
\maketitle

\begin{abstract}
Respiratory recordings provide several views of the same participant, but
multimodal gains can be overstated when recordings from one person cross data
partitions or when fusion is chosen on test results. We developed an audio-only
Coswara system using cough, breathing, and speech under fixed
participant-disjoint train, validation, and test partitions. Audio was quality
screened and represented by 10,140 numeric audio-derived candidates; feature ranking used
only training data and retained 800 candidates. Modality models and three
probability-fusion rules were evaluated using validation participants. The
selected speech branch achieved test AUROC 0.888, while the validation-selected
cough--speech uniform mean achieved AUROC 0.895 (95\% participant-bootstrap
CI 0.852--0.933), AUPRC 0.862, balanced accuracy 0.808, and F1 0.729 on 314
participants. The fusion estimate was 0.007 AUROC above speech, but paired
uncertainty for that difference was not estimated; the result does not
establish statistical superiority, external validity, or clinical validity.
\end{abstract}

\begin{IEEEkeywords}
COVID-19, respiratory audio, multimodal fusion, participant-disjoint evaluation,
acoustic features
\end{IEEEkeywords}

\section{Introduction}

Smartphone cough, breathing, and voice recordings have been studied as
low-cost inputs for COVID-19 risk assessment. Coswara made several prompted
respiratory sounds available for each participant and established a common
benchmark for this work \cite{bhattacharya2023coswara}. Subsequent studies
used transfer learning across cough, breathing, and speech
\cite{pahar2022transfer}, combined acoustic and symptom information
\cite{chetupalli2023multimodal}, and developed attention or transformer models
for respiratory spectrograms \cite{truong2024fair,aytekin2024hst}. These
results show that multiple sound types can be informative, but they do not make
fusion automatically beneficial: each added stream can introduce missingness,
noise, and another opportunity for model selection.

Evaluation design is equally important. Han \emph{et al.} showed that biased
cohort construction can materially change respiratory-audio performance
\cite{han2022realistic}. Coppock \emph{et al.} further demonstrated that
matching and longitudinal evaluation can reduce estimates obtained under a
random participant split \cite{coppock2024audio}. Feature selection is another
source of optimism; in the DiCOVA setting, Avila \emph{et al.} reported a
substantial difference between nested and non-nested selection
\cite{avila2021feature}. These findings motivate a source-development study in
which participant identities are isolated and the reported candidate ranking
is recomputed from training and validation evidence without using test metrics.

This paper asks a focused question: within a fixed participant-disjoint Coswara
evaluation, which fusion rule is selected without test-guided model choice,
and how does its held-out estimate compare with the selected single modalities?
The contribution is threefold. First, we combine
standardized openSMILE descriptor sets with a broad signal-summary bank and
perform modality-balanced, training-only feature ranking. Second, we compare
uniform, validation-weighted, and validation-fitted logistic fusion while
excluding test metrics from the reported selection rule. Third, we report the
corresponding test result with
participant-level bootstrap uncertainty. The study is
limited to source-domain model development; temporal and cross-dataset
transport are separate research questions.

\section{Materials and Methods}

\subsection{Cohort and Participant Partitioning}

The indexed Coswara release contained 2,746 participants, of whom 2,114 had
resolved active-positive or healthy/negative labels. After recording-quality
screening, 2,088 participants retained at least one eligible cough, breathing,
or speech recording and entered the modeling cohort; unresolved-label records
were not fitted or evaluated. A participant contributed up to two coughs, two
breathing recordings, two counting recordings, and three sustained vowels.
All recordings from one participant remained in one fixed partition. Table
\ref{tab:cohort} reports the partition counts; the final cough--speech test
analysis used the 314 participants for whom both selected modality predictions
were available. The initial 70/15/15 participant assignment used seed 42 and
stratified on label, age band, and gender where cell sizes permitted, with
fallback to label-only stratification.

\begin{table}[!b]
\caption{Quality-Passing Participant-Disjoint Coswara Partitions}
\label{tab:cohort}
\centering
\footnotesize
\begin{tabular}{lrrrr}
\toprule
Partition & Participants & Positive & Negative & Recordings \\
\midrule
Train      & 1,460 & 476 & 984 & 12,685 \\
Validation &   312 & 101 & 211 &  2,702 \\
Test       &   316 & 103 & 213 &  2,719 \\
\bottomrule
\end{tabular}
\end{table}

\subsection{Signal Processing and Acoustic Representation}

Recordings were decoded as mono waveforms and resampled to 16~kHz. A 35-dB
silence rule removed low-energy boundaries, frame-level root-mean-square
energy localized the active region, and peak normalization preceded a centered
crop or zero padding to 4~s for cough and 5~s for breathing and speech.
Only recordings flagged ``ok'' were modeled: minimum duration was 0.5~s for
cough and 1~s for breathing or speech, silence ratio could not exceed 0.70,
and clipping ratio could not exceed 0.01.

Each recording was represented by three complementary banks. The canonical
openSMILE ComParE~2016 and IS10 sets supplied 6,373 and 1,582 functionals over
energy, spectral, cepstral, and voicing descriptors
\cite{eyben2010opensmile,schuller2010is10,schuller2016compare}. A project bank
summarized 40 MFCC trajectories and their first and second differences, 64 mel
bands, chroma, spectral contrast, tonnetz, energy, zero-crossing, spectral
shape, and tempo measurements. Four preprocessing descriptors---event start,
event end, event duration, and active-audio ratio---were appended to each
openSMILE table. After merging, 10,140 numeric candidates were available:
6,377 in the ComParE-derived table, 1,586 in the IS10-derived table, and 2,177
project features.

Feature ranking was fitted only to the training partition. LightGBM gain for
feature $j$ was normalized within modality $m$ and averaged across modalities,
\begin{equation}
s_j=\frac{1}{|\mathcal{M}|}\sum_{m\in\mathcal{M}}
\frac{g_{jm}}{\max_k g_{km}},
\end{equation}
and the 800 highest-scoring feature names were frozen for validation and test.
This max-normalized rule prevents the modality with the largest raw gain scale
from determining the common bank. The final-validation analysis used the fixed
800-feature configuration; 500 and 1,200 were also examined during development,
so inference here is specific to the 800-feature setting. Within each classifier pipeline,
constant columns were removed and an ANOVA F-score filter retained 80\% of the
800 candidates for the boosted models or 60\% for the radial-basis SVC.

\subsection{Modality Models and Fusion}

Four classifier families were used in the final comparison: LightGBM,
XGBoost, CatBoost, and a sigmoid-calibrated radial-basis SVC
\cite{cortes1995svm,ke2017lightgbm,chen2016xgboost,prokhorenkova2018catboost}.
XGBoost used 800 trees of depth 3; LightGBM used 900 trees and 31 leaves;
CatBoost used 900 trees of depth 5. All three used learning rate 0.025 and the
boosted pipelines retained the top 80\% of candidates after removing constant
columns. The SVC used $C=2$, scaled inputs, the top 60\% of candidates, and
three-fold sigmoid calibration. All stochastic components used seed 42.
For the boosted models, SMOTE with three nearest neighbors was applied only to
training rows after feature filtering \cite{chawla2002smote}; validation and
test distributions were never resampled. Recording probabilities were first
averaged within participant and modality. For each modality, validation AUROC
(with AUPRC as a tie breaker) selected either one classifier or the mean of the
four validation-ranked classifiers. This selected a four-model mean for cough
and breathing and LightGBM for speech.

Let $p_m(x)$ denote the selected probability for modality $m$ and participant
$x$. Uniform fusion averaged the available modality probabilities. A second
rule used validation AUPRC:
\begin{equation}
\begin{aligned}
r_m&=\max(\operatorname{AUPRC}_m-0.5,0.01),\\
w_m&=r_m\big/\textstyle\sum_j r_j,\\
p_w(x)&=\textstyle\sum_m w_m p_m(x).
\end{aligned}
\end{equation}
For cough and speech this produced weights 0.459 and 0.541, respectively. This
0.5 offset was an exploratory heuristic, not a prevalence-derived no-skill
baseline. The third rule fitted class-balanced L2 logistic regression directly
to aligned validation probabilities,
\begin{equation}
p_s(x)=\sigma\!\left(\beta_0+\sum_m\beta_mp_m(x)\right).
\end{equation}
Its coefficients were learned from validation participants rather than
hard-coded. Because those same participants generated its validation score,
the stack and heuristic weighting were treated as exploratory. The primary
rule was fixed uniform averaging; validation AUROC and AUPRC selected only its
modality combination. A balanced-accuracy threshold was selected on validation
data. This validation-first rule was applied retrospectively to saved candidate
rows: code determined the candidate key from validation columns and then
retrieved its corresponding test row without ranking on test metrics.

\begin{figure*}[!t]
\centering
\includegraphics[width=\textwidth]{conference_design_and_results.pdf}
\caption{Participant-isolated model development and its test result. (a) Only
training participants informed feature ranking and model fitting; validation
participants determined modality models, fusion, and threshold, while test
metrics were excluded from selection. (b) Participant-level AUROC and AUPRC point estimates for the
selected single modalities and validation-selected cough--speech mean.}
\label{fig:design-results}
\end{figure*}

\subsection{Evaluation and Uncertainty}

The primary endpoint was participant-level AUROC; AUPRC was reported because
it reflects the positive-class ranking more directly. At the frozen validation
threshold we also calculated balanced accuracy and F1. A nonparametric
participant bootstrap resampled test participants 1,000 times to obtain a 95\%
confidence interval for the selected AUROC. This interval quantifies sampling
uncertainty within the fixed test cohort; it does not estimate variability
from model refitting or from drawing a new cohort.

\section{Results}

\subsection{Validation-Only Selection}

Among the four-family final candidate bank, validation AUROC selected the
four-model mean for breathing (0.783), the four-model mean for cough (0.807),
and LightGBM for speech (0.843). Among the predefined uniform-mean
combinations, cough plus speech ranked first (AUROC 0.842; AUPRC 0.800),
slightly above the three-modality mean (AUROC 0.841). This validation-only rule
identified the primary configuration, whose corresponding test row was then
retrieved. The fitted stack's
in-sample validation AUROC of 0.844 was not eligible for primary selection.

\subsection{Participant-Disjoint Test Performance}

Table \ref{tab:results} gives the controlled test comparison. Speech was the
strongest selected single modality with AUROC 0.888 and AUPRC 0.833. The
primary cough--speech mean reached AUROC 0.895 (95\% CI 0.852--0.933), AUPRC
0.862, balanced accuracy 0.808, and F1 0.729. Relative to speech alone, the
descriptive differences were +0.007 AUROC, +0.028 AUPRC, approximately zero
balanced accuracy, and $-0.011$ F1. The heuristic weighted mean and fitted
stack reached AUROC 0.896 and 0.897, respectively, but remained exploratory.

\begin{table}[!t]
\caption{Participant-Level Test Results}
\label{tab:results}
\centering
\footnotesize
\setlength{\tabcolsep}{3.2pt}
\begin{tabular}{lrrrrr}
\toprule
System & $n$ & AUROC & AUPRC & Bal. Acc. & F1 \\
\midrule
Breathing, top-4 mean & 312 & 0.828 & 0.734 & 0.745 & 0.654 \\
Cough, top-4 mean     & 316 & 0.862 & 0.804 & 0.785 & 0.705 \\
Speech, LightGBM      & 314 & 0.888 & 0.833 & 0.809 & 0.740 \\
\midrule
\textbf{Cough+speech mean} & \textbf{314} & \textbf{0.895} & \textbf{0.862} & \textbf{0.808} & \textbf{0.729} \\
Cough+speech weighted$^*$ & 314 & 0.896 & 0.863 & 0.820 & 0.744 \\
Cough+speech stack$^*$    & 314 & 0.897 & 0.863 & 0.825 & 0.752 \\
\bottomrule
\end{tabular}
\par\vspace{1pt}\scriptsize $^*$Exploratory; not eligible for primary selection.
\end{table}

\section{Discussion}

The validation procedure selected cough plus speech for the fixed uniform-mean
rule, and its held-out AUROC was numerically 0.007 above the selected speech
branch. AUPRC increased, but balanced accuracy was unchanged and F1 decreased.
Paired uncertainty was unavailable, so this is not evidence of statistical
superiority. Breathing was available to the search but absent from the primary
configuration. The result supports evaluating multimodality selectively rather
than assuming that every additional sound stream helps.

\begin{table}[H]
\caption{Nearest Internal Coswara Context}
\label{tab:context}
\centering
\footnotesize
\setlength{\tabcolsep}{3.1pt}
\begin{tabular}{lll}
\toprule
Study & Inputs & AUROC \\
\midrule
Coswara \cite{bhattacharya2023coswara} & 9 sounds + symptoms & 0.915 \\
Chetupalli \cite{chetupalli2023multimodal} & Audio-only fusion & 0.880 \\
FAIR \cite{truong2024fair} & 7 sound instances & $0.866\pm0.012$ \\
This work & Cough + speech mean & 0.895 \\
\bottomrule
\end{tabular}
\end{table}

Table \ref{tab:context} places the AUROC in the nearest internal Coswara
context. Chetupalli \emph{et al.} reported 0.880 for audio-only fusion under a
subject-disjoint protocol \cite{chetupalli2023multimodal}; FAIR reported
$0.866\pm0.012$ using seven sound instances and frozen encoders
\cite{truong2024fair}. The Coswara dataset paper reported 0.915 using nine
sound tasks together with symptoms \cite{bhattacharya2023coswara}. These are
not leaderboard-equivalent values: cohort snapshot, label filtering, inputs,
and partition construction differ. They show that the present audio-only
estimate is competitive within internal Coswara studies, not that it is
statistically superior to them.

This study has five main limitations. First, it is a retrospective analysis of
one crowdsourced resource; the dataset-provided COVID status was not
independently adjudicated by the present authors. Second, validation data served
several selection roles, including modality model, modality combination, and
threshold choice. The validation-first primary ranking was reconstructed
retrospectively rather than preregistered, and the 800-feature budget was one
of three explored settings;
a nested outer evaluation would estimate the complete procedure more
conservatively. Third, event timing and active-audio descriptors can encode
recording protocol, and repeated recordings were fitted as separate rows; no
ablation isolated these effects. Fourth, only 314 test participants had the
complete selected cough--speech predictions. The bootstrap interval describes
the primary AUROC, not its paired difference from speech, so the 0.007
difference is descriptive. Fifth, calendar drift, independent datasets, and
prospective clinical use are outside this experiment and require separate
protocols.

\section{Conclusion}

Training-only acoustic ranking followed by validation-guided probability
fusion selected the cough--speech mean as the primary configuration. Its
participant-disjoint test AUROC was 0.895, compared with 0.888 for the selected
speech branch. Without paired uncertainty, this numerical difference is not a
superiority claim. The evidence supports careful
multimodal model development rather than diagnostic readiness. Any extension to a
new architecture, time period, or dataset should preserve the same separation
between fitting, selection, and final evaluation.

\section*{Ethics and Data Availability}

This work is a secondary analysis of the public Coswara resource; no new
participants were recruited. Participant information, consent, and collection
procedures are described in the source publication
\cite{bhattacharya2023coswara}. Raw recordings are not redistributed. Analysis
code, frozen configurations, and aggregate result tables will accompany the
archived project release.

\IEEEtriggeratref{9}
\bibliographystyle{IEEEtran}
\bibliography{references}

\end{document}
